% Hafta 9 — MATE: saldırgan cihazın sahibi
% Derleme: py -3.12 tools/latex/derle.py docs/week-9/assets/h09-03-mate.tex
\documentclass[tikz,border=8pt]{standalone}
\usepackage{cen429-cizim}
\begin{document}
\begin{tikzpicture}[node distance=10mm and 10mm]
  \node[baslik] (b) at (0,0) {MATE: saldırgan cihazın sahibi};
  \node[altbaslik, text width=170mm] at (b.south west) {mobil ödeme, DRM, lisanslama: kod kullanıcının elinde çalışır};

  \node[kutu=78mm, below=10mm of b.south west, anchor=north west] (sol) {\kb{Ağ saldırganı}\\yalnız mesajları görür\\TLS onu durdurur\\[2mm]klasik güvenlik dersleri\\bunu varsayar};
  \node[kotukutu=78mm, right=10mm of sol] (sag) {\kb{MATE (uçtaki insan)}\\ikiliyi açar, belleği okur\\debugger bağlar, yamalar\\istediği kadar dener\\[2mm]kripto TEK BAŞINA yetmez};
  \draw[draw=cizgi, dashed] ($(sol.north east)!0.5!(sag.north west)$) -- ($(sol.south east)!0.5!(sag.south west)$);

  \node[dip, text width=170mm, below=10mm of sag.south -| sol, anchor=north west] {%
    Gizleme burada devreye girer: kırılamazlık değil, MALİYET yükseltme.};
\end{tikzpicture}
\end{document}
